\documentclass[10pt,a4paper]{amsart}
\usepackage[margin=19mm]{geometry}
\usepackage{amsmath,amssymb,mathtools}
\usepackage[scr=rsfs,cal=euler]{mathalfa}
\usepackage{xfrac}
\usepackage[unicode,hidelinks,bookmarks=false]{hyperref}
\newcommand{\ie}{i.e.\ }
\newcommand{\cf}{cf.\ }
\newcommand{\resp}{respectively\ }
\newcommand{\Subsection}{\subsection}
\providecommand{\nobreakdash}{\nobreak\hbox{-}}
\setlength{\emergencystretch}{2em}
\multlinegap0pt
\usepackage{amssymb}
\usepackage{dsfont}
\usepackage{xcolor}
\usepackage[shortlabels]{enumitem}
\newtheorem{prop}{Proposition}
\newtheorem{thm}{Theorem}
\renewcommand{\Pr}{\operatorname{pr}}
\title{Consistency of variational approximations\\ under bounded Kullback--Leibler divergence}
\author{Hien Duy Nguyen, Jacob Westerhout, Thomas Guilmeau, Julyan Arbel}
\date{Source: arXiv:2606.13230v1 (CC BY 4.0)}
\begin{document}
\maketitle

\noindent\emph{Source and licence.} Consistency of variational approximations\\ under bounded Kullback--Leibler divergence. Version arXiv:2606.13230v1 (CC BY 4.0), distributed by arXiv under the Creative Commons Attribution 4.0 International licence, http://creativecommons.org/licenses/by/4.0/, as recorded in the arXiv licence metadata for this exact version and shown on the abstract page https://arxiv.org/abs/2606.13230v1. Full licence notice: \texttt{licenses/arxiv-2606.13230-CC-BY-4.0.txt}.\par\smallskip

\noindent\textit{Editorial scope.} Counted unit: the article's thm thm:gp-envelope (source0.tex lines 431--465). The article's complete original proof of it is delivered with it (source0.tex lines 708--796, one \texttt{proof} environment of 89 lines, which the article places away from the statement). No mathematics is written, repaired or re-derived: the statement and the proof are the article's own lines, in the article's own notation, and no retained line is substituted or re-flowed.  Retained uncounted as the source-local context the proof needs: lines 384--394.  Nothing was omitted from any retained span, so the package carries no omission record.  No retained line contains a citation, so the wrapper carries no bibliography and no bibliography entry.  No retained line contains a figure environment, an image-inclusion command or drawing code, so no image is delivered and the package declares no asset dependency.  Editorial formatting only, in addition: an AMS \texttt{amsart} preamble that restates the article's own declarations and macros used by the retained lines, namely the environments \texttt{prop}, \texttt{thm} on separate counters, together with the standard packages those lines need.  The retained environments print in this document as Proposition 1, Theorem 1 in order of appearance, whereas the article prints the same items with the numbering of its own full document.  The complete native source of the article as distributed in the arXiv e-print for this exact version is bundled unchanged as \texttt{sources/202688/source0.tex}.
\begin{prop} \label{prop:quadratic-envelope} Let $\tilde\mu_{n}$ be Borel
probability measures on $\mathbb{R}^{p}$ with strictly positive Lebesgue
densities $\tilde p_{n}$. Write $g_{n}=\log \tilde p_{n}$. Assume that there exist
constants $L,M,\alpha,r>0$ and $n_{0}\in\mathbb{N}$ such that, for
every $n\ge n_{0}$, $\mathrm{(B1)}$  $\nabla g_{n}$ is $L$-Lipschitz on $\mathbb{R}^{p}$, $\mathrm{(B2)}$  $\|\nabla g_{n}(0)\|\le M$, $\mathrm{(B3)}$ $\tilde\mu_{n}(B(0,r))\ge\alpha$.
Then there exist constants $c,C>0$, depending only on $L,M,\alpha,r$
and $p$, such that
\[
\tilde p_{n}(z)\ge c\exp\left(-\frac{C}{2}\|z\|^{2}\right)\qquad\text{for all }z\in\mathbb{R}^{p}\text{ and all }n\ge n_{0}.
\]
\end{prop}

\begin{thm} \label{thm:gp-envelope} Assume:
\begin{enumerate}
\item[$\mathrm{(C1)}$] For every $x\in\mathbb{X}$, the map $\theta\mapsto\ell(\theta,x)$
is twice continuously differentiable, and there exists a measurable
$H:\mathbb{X}\to[0,\infty)$ such that
\[
\sup_{\theta\in\mathbb{R}^{p}}\|\nabla_{\theta}^{2}\ell(\theta,x)\|\le H(x),\qquad\mathrm{E}[H(X_{1})]<\infty.
\]
\item[$\mathrm{(C2)}$] $\pi$ is strictly positive and twice continuously differentiable
on $\mathbb{R}^{p}$, and
\[
M_{\pi}=\sup_{\theta\in\mathbb{R}^{p}}\|\nabla_{\theta}^{2}\log\pi(\theta)\|<\infty.
\]
\item[$\mathrm{(C3)}$] $\bar{\eta}=\sup_{n\ge1}\eta_{n}<\infty$.
\item[$\mathrm{(C4)}$] There exist almost sure events $\Omega_{\mathrm{loss}},\Omega_{\nabla}\in\mathfrak{A}$
such that, for every $\omega\in\Omega_{\mathrm{loss}}$ and every
$n\ge1$, $\hat{\theta}_{n}(\omega)$ is a global minimiser of $\theta\mapsto\sum_{i=1}^{n}\ell(\theta,X_{i}(\omega))$,
and, for every $\omega\in\Omega_{\nabla}$,
\[
n^{-1/2}\left\|\nabla_{\theta}\log\pi\left(\hat{\theta}_{n}(\omega)\right)\right\|\to0.
\]
\item[$\mathrm{(C5)}$] There exist a Borel probability measure $\mu_{\infty}$ on $\mathbb{R}^{p}$,
a radius $r>0$, and an event $\Omega_{\mathrm{w}}\in\mathfrak{A}$
with $\Pr(\Omega_{\mathrm{w}})=1$ such that, for every $\omega\in\Omega_{\mathrm{w}}$,
\[
\tilde\mu_n(\omega)\rightsquigarrow\mu_{\infty}\qquad\text{and}\qquad\mu_{\infty}(B(0,r))>0.
\]
\end{enumerate}
Then there exist constants $c,C>0$ and an event $\Omega_{0}\in\mathfrak{A}$
with $\Pr(\Omega_{0})=1$ such that, for every $\omega\in\Omega_{0}$,
there exists $n_{0}(\omega)\in\mathbb{N}$ for which
\begin{equation}
\tilde p_{n}(\omega,z)\ge c\exp\left(-\frac{C}{2}\|z\|^{2}\right)\qquad\text{for all }z\in\mathbb{R}^{p}\text{ and all }n\ge n_{0}(\omega).\label{eq:Gen_Bayes_Lower_bound}
\end{equation}
\end{thm}

\begin{proof}[of Theorem~\ref{thm:gp-envelope}] By the strong
law of large numbers,
\[
\Omega_{H}=\left\{\omega\in\Omega: \frac{1}{n}\sum_{i=1}^{n}H(X_{i}(\omega))\to\mathrm{E}[H(X_{1})]\right\}
\]
has probability one. Let
\[
\Omega_{0}=\Omega_{H}\cap\Omega_{\mathrm{loss}}\cap\Omega_{\nabla}\cap\Omega_{\mathrm{w}},
\]
so that $\Pr(\Omega_{0})=1$. Fix $\omega\in\Omega_{0}$ for the remainder
of the proof and suppress the $\omega$-dependence when no confusion
arises.

For each $n$, define
\[
g_{n}(z)=\log \tilde p_{n}(z),\qquad z\in\mathbb{R}^{p}.
\]
Since
\[
g_{n}(z)=-\eta_{n}\sum_{i=1}^{n}\ell\!\left(\hat{\theta}_{n}+n^{-1/2}z,X_{i}\right)+\log\pi\!\left(\hat{\theta}_{n}+n^{-1/2}z\right)-\log\mathcal{Z}_{n}-\frac{p}{2}\log n,
\]
for a normalising constant $\mathcal{Z}_{n}$ independent of $z$,
assumptions {(C1)} and {(C2)} imply that $g_{n}$ is twice
continuously differentiable and
\[
\nabla_{z}^{2}g_{n}(z)=\frac{1}{n}\left\{ -\eta_{n}\sum_{i=1}^{n}\nabla_{\theta}^{2}\ell\!\left(\hat{\theta}_{n}+n^{-1/2}z,X_{i}\right)+\nabla_{\theta}^{2}\log\pi\!\left(\hat{\theta}_{n}+n^{-1/2}z\right)\right\} .
\]
Hence
\[
\sup_{z\in\mathbb{R}^{p}}\|\nabla_{z}^{2}g_{n}(z)\|\le\eta_{n}\frac{1}{n}\sum_{i=1}^{n}H(X_{i})+\frac{M_{\pi}}{n}.
\]
Since $\omega\in\Omega_{H}$, there exists $n_{1}(\omega)\in\mathbb{N}$
such that, with $L=\bar{\eta}\left(\mathrm{E}[H(X_{1})]+1\right)+M_{\pi}$,
\[
\sup_{z\in\mathbb{R}^{p}}\|\nabla_{z}^{2}g_{n}(z)\|\le L\qquad\text{for all }n\ge n_{1}(\omega).
\]
In particular, $\nabla g_{n}$ is $L$-Lipschitz on $\mathbb{R}^{p}$
for all $n\ge n_{1}(\omega)$.

Next,
\[
\nabla_{z}g_{n}(0)=n^{-1/2}\left\{ -\eta_{n}\sum_{i=1}^{n}\nabla_{\theta}\ell(\hat{\theta}_{n},X_{i})+\nabla_{\theta}\log\pi(\hat{\theta}_{n})\right\} .
\]
Because $\hat{\theta}_{n}$ is a global minimiser of $\theta\mapsto\sum_{i=1}^{n}\ell(\theta,X_{i})$
and this criterion is continuously differentiable on $\mathbb{R}^{p}$,
we have
\[
\sum_{i=1}^{n}\nabla_{\theta}\ell(\hat{\theta}_{n},X_{i})=0.
\]
Therefore
\[
\nabla_{z}g_{n}(0)=n^{-1/2}\nabla_{\theta}\log\pi(\hat{\theta}_{n}).
\]
Since $\omega\in\Omega_{\nabla}$, there exists $n_{2}(\omega)\in\mathbb{N}$
such that
\[
\|\nabla_{z}g_{n}(0)\|\le1\qquad\text{for all }n\ge n_{2}(\omega).
\]

Since $\omega\in\Omega_{\mathrm{w}}$, we have
\[
\tilde\mu_n\rightsquigarrow\mu_{\infty}\qquad\text{and}\qquad\mu_{\infty}(B(0,r))>0.
\]
The Portmanteau theorem gives
\[
\liminf_{n\to\infty}\tilde\mu_n(B(0,r))\ge\mu_{\infty}(B(0,r))>0.
\]
Hence, with
\[
\alpha=\frac{1}{2}\mu_{\infty}(B(0,r))>0,
\]
there exists $n_{3}(\omega)\in\mathbb{N}$ such that
\[
\tilde\mu_n(B(0,r))\ge\alpha\qquad\text{for all }n\ge n_{3}(\omega).
\]

Applying Proposition~\ref{prop:quadratic-envelope} to the deterministic
sequence $(\tilde\mu_n(\omega))_{n\ge1}$ with the constants $L$, $M=1$,
$\alpha$, and $r$, we obtain constants $c,C>0$, depending only
on $\bar{\eta},~\mathrm{E}[H(X_{1})],~M_{\pi},~r,~\mu_{\infty}(B(0,r)),~p$,
and an index
\[
n_{0}(\omega)=\max\{n_{1}(\omega),n_{2}(\omega),n_{3}(\omega)\}
\]
such that
\[
\tilde p_{n}(\omega,z)\ge c\exp\left(-\frac{C}{2}\|z\|^{2}\right)\qquad\text{for all }z\in\mathbb{R}^{p}\text{ and all }n\ge n_{0}(\omega).
\]
\end{proof}

\end{document}
